\documentclass[../../../include/open-logic-section]{subfiles}

\begin{document}
\olfileid{sth}{ordinals}{zfm}	
\olsection{$\ZFminus$: يو مهم پړاو}

دا پوښتنه چې تعويض څنګه توجيه کړو، که ئې توجيه کېدلے شي، اسانه نۀ ده.
نو دا به د \olref[replacement][]{chap} دپاره پرېږدو. خو د تعويض په اضافه
کولو سره يو بل مهم پړاو ته رسېدلي يو. اوس زموږ سره د $\ZFminus$ د
نظريې ټول اړين بديهي اصلونه شته. په تفصيل سره:

\begin{defn}
نظريه $\ZFminus$ دا بديهي اصلونه لري: د غړيتوب له مخې برابري، اتحاد،
جوړې، د فرعي سټونو سټونه، نامتناهيت، او د بېلولو او تعويض د شېماګانو
ټولې بېلګې. په بله وينا، $\ZFminus$ په $\Zminus$ کښې تعويض اضافه کوي.
\end{defn}
\noindent
دا د \emph{زرمېلو--فرېنکل} د سټونو د نظريې نوم دے، چې يو څه ترې
\emph{کم} دي؛ هغې خبرې ته به وروسته راشو. فرېنکل ته دا وياړ ځکه ورکول
کېږي چې په \citeyear{Fraenkel1922} کښې د تعويض فورمولبندي ده ته منسوبه
ده، که څه هم لومړۍ دقيقه فورمولبندي د \citet{Skolem1922} وه.

\end{document}
